\documentclass[conference]{IEEEtran}
% =====================================================================
%  中文译版：Run, Read Back, Refine: Observation Fidelity in
%            Execution-Grounded EDA Agents
%  原文：arXiv:2611.19056v2，第一作者 Yang Liu
%  中文版由原作者授权发布。本译版仅供学习与研究，非商业用途。
%
%  构建方式（中文需 XeLaTeX）：
%      xelatex main && bibtex main && xelatex main && xelatex main
% =====================================================================

% ---------------- packages (IEEE template set + what this paper needs) --
\usepackage{cite}
\usepackage{amsmath,amssymb,amsfonts}
\usepackage{algorithmic}
\usepackage{graphicx}
\usepackage{textcomp}
\usepackage{xcolor}
\usepackage{booktabs}
\usepackage{url}

% ---------------- 中文支持 -------------------------------------------
% xeCJK 负责中文断行与字体；IEEEtran 保持 Times 正文体例
\usepackage{xeCJK}
% 文泉驿正黑：TTF 轮廓，标点与汉字覆盖完整（Droid Sans Fallback 缺弯引号与破折号）
\setCJKmainfont{WenQuanYi Zen Hei}
\setCJKsansfont{WenQuanYi Zen Hei}
\setCJKmonofont{WenQuanYi Zen Hei}
\xeCJKsetup{CJKecglue={\hskip 0.08em plus 0.04em minus 0.02em}}
% 中文段落：首行缩进两字，取消英文段间距
\setlength{\parindent}{2em}
\setlength{\parskip}{0pt}

\setlength{\textfloatsep}{8pt plus 2pt minus 2pt}
\setlength{\floatsep}{6pt plus 2pt minus 2pt}
\setlength{\intextsep}{6pt plus 2pt minus 2pt}
\setlength{\abovecaptionskip}{3pt}
\setlength{\belowcaptionskip}{0pt}

\begin{document}

% ---------------- title ------------------------------------------------
\title{运行、回读、精炼：\\执行接地 EDA 智能体中的观测保真度}

% ---------------- 作者 --------------------------------------------------
\author{%
\IEEEauthorblockN{Yang Liu, Shiwei Hou, Xiyuan Chen, Sen Yuan, Shao You, Feifan Chen, Wencheng Li, Shuyang Hu,\\
Emma Xia, Xiaojing Lu, Wenchuan Zhang, Hao Wang, Fan Xu, Yanfeng Li, Hongwen Li, Yiming Zhu}\\
\IEEEauthorblockA{\textit{长鑫存储技术有限公司}，中国；\quad 邮箱：\{paul.liu, yanfeng.li, howie.li\}@cxmt.com}
}

\maketitle

% ---------------- 行距（中文需拉开；必须放 \maketitle 之后）-------------
% ⚠️ 不要在导言区改 \linespread / \normalsize / \baselineskip —— IEEEtran 双栏浮动体
%    取样依赖它们的原始值，改了会导致跨栏断句、两栏交错（详见 XIANLUN_TRANSLATE_SKILL.md §3.3）
\renewcommand{\baselinestretch}{1.35}\selectfont

% ---------------------------------------------------------------- 摘要
% 译版授权说明只留在文件头注释里，不排进正文（见 SKILL §3.4）

% ---------------- abstract + index terms -------------------------------
\input{DAC2027/0_abstract}

% ---------------- body -------------------------------------------------
\input{DAC2027/1_intro}
\input{DAC2027/2_background}
\input{DAC2027/3_design}
\input{DAC2027/4_grounding}
\input{DAC2027/5_bench}
\input{DAC2027/6_exp}
\input{DAC2027/7_discussion}
\input{DAC2027/8_conclusion}

% ---------------- references -------------------------------------------
\clearpage
\bibliographystyle{IEEEtran}
\bibliography{reference}

\end{document}
